\documentclass[12pt]{article}
\usepackage[T1]{fontenc}
\usepackage[utf8]{inputenc}
\usepackage{lmodern}
\usepackage{textcomp}
\usepackage{enumitem}
\usepackage{geometry}
\geometry{margin=1in}
\begin{document}
\begin{center}
Answer Key - Biology - Graduate Level Assessment 13 \
\small (Answers are ordered exactly as in the question paper.)
\end{center}

\section*{Multiple Choice}
\begin{enumerate}[label=\arabic*.]
\item \textbf{Answer:} D
\\ \textit{Options:} A) Glomerular capillary pressure and blood volume; B) Diameter of the efferent arterioles and plasma protein concentration; C) Renal blood flow and the concentration of sodium in the distal tubule; D) Arterial blood pressure and diameter of the afferent arterioles; E) Kidney size and blood urea nitrogen levels
\\ \textit{Note:} The glomerular filtration rate (GFR) is primarily influenced by changes in glomerular capillary pressure, which is regulated by arterial blood pressure and the diameter of the afferent arterioles, as these factors directly affect the blood flow and pressure entering the glomerulus.
\item \textbf{Answer:} A
\\ \textit{Options:} A) Oxygen is used to maintain the temperature of cells.; B) Oxygen is used by cells to produce alcohol.; C) Oxygen is used to produce carbon dioxide in cells.; D) Oxygen is solely responsible for stimulating brain activity.; E) Oxygen is used to synthesize proteins in cells.
\\ \textit{Note:} Oxygen is critical for cellular respiration, which produces the energy necessary for maintaining metabolic processes that regulate cell temperature and overall homeostasis.
\item \textbf{Answer:} E
\\ \textit{Options:} A) Epilepsy is the result of a chromosomal aneuploidy, similar to Down syndrome.; B) Epilepsy is always inherited in an X-linked recessive pattern.; C) Epilepsy is a dominant trait.; D) Epilepsy is caused by an autosomal dominant trait.; E) Epilepsy is due to an autosomal recessive trait.
\\ \textit{Note:} The correct answer identifies epilepsy as an autosomal recessive trait, which is supported by genetic studies showing that certain types of epilepsy can be inherited in this manner, highlighting the role of genes in the etiology of the disorder.
\item \textbf{Answer:} A
\\ \textit{Options:} A) that humans evolved from apes.; B) that evolution occurs.; C) that evolution occurs through a process of gradual change.; D) that the Earth is older than a few thousand years.; E) a mechanism for how evolution occurs.
\\ \textit{Note:} The correct answer is based on the historical context of Darwin's theory of evolution through natural selection, which fundamentally challenged the prevailing views of human origins by suggesting that humans share a common ancestor with apes, rather than being a separate creation.
\item \textbf{Answer:} A
\\ \textit{Options:} A) Mimetic coloration; B) Aposomatic coloration; C) Deceptive markings; D) Batesian mimicry; E) Flash coloration
\\ \textit{Note:} Mimetic coloration is the adaptive strategy where organisms evolve colors and patterns that resemble their environment, enhancing their ability to avoid predation by blending in with surrounding elements.
\end{enumerate}

\section*{True / False}
\begin{enumerate}[label=\arabic*.]
\setcounter{enumi}{5}
\item \textbf{Answer:} False
\\ \textit{Options:} A) True; B) False
\\ \textit{Note:} This study showed that FDG is not a good surrogate tracer for tumor hypoxia under either ambient or hypoxic conditions. Only specific hypoxia tracers should be used to measure tumor hypoxia.
\item \textbf{Answer:} True
\\ \textit{Options:} A) True; B) False
\\ \textit{Note:} The prevalence of vagino-rectal colonization by Streptococcus group B in the pregnant women from Melilla is within the national estimated figures, however it is different if they are from Muslim or Christian culture, being higher in the Muslim population. On one hand both prevalences are within the national statistics, and on the other hand it is observed that there is not any difference according to age.
\item \textbf{Answer:} True
\\ \textit{Options:} A) True; B) False
\\ \textit{Note:} GFRUPs procedure was applicable in most cases. The main difficulties were anticipating the correct date for the meeting and involving nurses in the procedure. Children for whom the procedure was interrupted because of clinical improvement and who survived in poor condition without a formal decision pointed out the need for medical criteria for questioning, which should systematically lead to a formal decision-making process.
\item \textbf{Answer:} True
\\ \textit{Options:} A) True; B) False
\\ \textit{Note:} According to the results from this meta-analysis, the risk of grade 3 and 4 CVAEs in patients who were receiving AIs was higher compared with the risk in patients who were receiving tamoxifen, and the difference reached statistical significance. However, the AD was relatively low, and from 160 to 180 patients had to be treated to produce 1 event.
\item \textbf{Answer:} True
\\ \textit{Options:} A) True; B) False
\\ \textit{Note:} The current literature suggests that dexamethasone can be used as an effective alternative to prednisone in the treatment of mild to moderate acute asthma exacerbations in children, with the added benefits of improved compliance, palatability, and cost. However, more research is needed to examine the role of dexamethasone in hospitalized children.
\end{enumerate}

\section*{Long Form Response}
\begin{enumerate}[label=\arabic*.]
\setcounter{enumi}{10}
\item \textbf{Answer:} Nerve gases that inhibit cholinesterase exert profound physiological effects on the human body by disrupting the normal degradation of the neurotransmitter acetylcholine at synapses. Normally, cholinesterase breaks down acetylcholine, allowing for the precise termination of nerve impulses. Inhibition of this enzyme leads to an accumulation of acetylcholine, resulting in continuous stimulation of muscles and glands, which can manifest as tremors, spasms, and potentially fatal respiratory failure due to paralysis of the diaphragm. The overstimulation of the nervous system can also lead to systemic symptoms such as bronchoconstriction, seizures, and loss of consciousness. Ultimately, the consequences of exposure to such nerve agents can be severe, culminating in death if not promptly treated with antidotes such as atropine and pralidoxime.
\\ \textit{Note:} correct because it accurately describes how the inhibition of cholinesterase leads to an accumulation of acetylcholine, causing persistent stimulation of muscles and glands, which results in severe physiological effects, including tremors, spasms, and respiratory failure.
\item \textbf{Answer:} A newborn can contract syphilis primarily through vertical transmission from an infected mother during pregnancy, particularly across the placenta. This can occur at any stage of gestation, with the risk of transmission increasing as the pregnancy progresses. Additionally, although rare, infants may also be exposed to syphilis through contact with infected lesions during delivery. The implications for maternal health are significant, as untreated syphilis can lead to severe complications such as miscarriage, stillbirth, or preterm birth, highlighting the importance of early detection and treatment in pregnant women. Furthermore, maternal syphilis can result in long-term health issues for the child, necessitating comprehensive prenatal care.
\\ \textit{Note:} correct because it accurately identifies vertical transmission as the primary mechanism through which a newborn contracts syphilis from an infected mother during pregnancy, acknowledges the timing of transmission risk, and includes the rare possibility of transmission through contact during delivery, while also addressing the serious implications of untreated syphilis for maternal health.
\end{enumerate}

\vfill
\noindent\textit{End of Answer Key}
\end{document}